\section{论文概览}

\subsection{核心定位}

FORGE 是一个面向制造场景的细粒度多模态评测基准。它关注的不是“模型能不能粗略看出这是螺丝或螺母”，而是通用多模态大语言模型（MLLM）能否在真实制造任务中完成从感知、细粒度语义理解到规则推理和决策判断的完整认知链条。

这篇论文的核心问题可以概括为：当前主流通用 MLLM 面对真实制造数据时，到底缺的是视觉定位能力，还是制造领域知识、细粒度形态理解和装配规则推理能力。实验结论倾向于后者：模型往往不是完全“看不见”，而是“不懂制造语义”和“不懂细粒度形态差异”。

\subsection{一句话总结}

FORGE 将真实 2D 图像、3D 点云、型号级语义标注和三类制造任务组织成一个标准化 benchmark，用来评估通用 MLLM 在制造场景下的细粒度识别、缺陷分析、装配验证和领域推理能力，并进一步证明这些标注也可以作为有效的监督微调资源。

\subsection{Benchmark 的含义}

benchmark 不是模型本身，而是一套标准化考试系统。它通常包含任务定义、题库数据、标准答案、评测协议和评分指标。dataset 更偏“原始数据资源”，benchmark 则进一步规定“用这些数据怎么出题、怎么答题、怎么评分、怎么比较模型”。

FORGE 既提供数据，也提供评测体系。它把真实制造图像和点云组织成工件验证、结构表面检测、装配验证三类任务，并采用统一的多项选择题协议与 exact-match accuracy 评分。因此，FORGE 不是单纯的数据集，而是一个围绕制造认知能力构建的评测框架。

\subsection{为什么是通用 MLLM 的认知能力评测}

论文关注的是通用 MLLM，而不是专门 3D 模型或传统工业视觉模型。通用 MLLM 指 GPT、Gemini、Claude、Qwen-VL、InternVL、Llama-Vision 等面向广泛任务的基础多模态模型。它们通常具有较强的语言推理、泛化和交互能力，因此理论上适合自动化制造中的高层判断与决策。

这里的“认知能力”不是哲学意义上的 cognition，而是工程上的多阶段能力：识别对象、理解规格、比较差异、应用制造规则、最终做出判断。普通视觉 benchmark 常问“这是什么”；FORGE 更进一步问“它具体是什么型号、是否满足装配要求、缺陷属于哪一类、为什么这个部件不合格”。这就是它从视觉评测上升到制造认知评测的原因。

\subsection{主要贡献}

\begin{enumerate}
    \item 构建真实制造场景中的多模态 benchmark，同时包含 2D 图像和 3D 点云。
    \item 引入型号级细粒度语义标注，覆盖工件类别、型号、缺陷类型、装配关系和制造约束。
    \item 设计三类贴近制造现场的任务：工件验证、结构表面检测和装配验证。
    \item 系统评估 18 个开源和闭源 MLLM，揭示当前模型的主要瓶颈不是 visual grounding，而是制造领域知识和 morphology understanding。
    \item 展示 FORGE 不只是静态评测集，也可作为 SFT 训练资源，帮助小模型获得可迁移的制造推理能力。
\end{enumerate}
